\section{MPE Implicit Solvent Model}

\emph{This is an experimental feature which is still in development. Do not rely on properties calculated by this method!}

The simulation of a solvent in a quantum mechanical calculation can, in principle, be done in two ways. One way is to include explicit solvent molecules in the calculation. This straightforward approach usually requires molecular dynamics (MD) simulations in order to yield thermodynamically meaningful observables as e.g.~solvation energies.

The second way is to average the effect of the solvent and treat it as a continuum which responds to the electrostatic potential created by the solute, i.e.~the entity that is to be solvated. There are several flavors to this comparatively inexpensive approximation, e.g.~the polarizable continuum model (PCM) by Mennucci and coworkers \cite{Mennucci02} or the conductor like screening model (COSMO) by Klamt and Sch\"u\"urmann \cite{Klamt93} to name two of the most popular ones.

The method presented here called multipole expansion (MPE) implicit solvent model is somewhat similar to the dielectric PCM model. It uses the multipole expansion of the electronic density to calculate the so-called reaction field, i.e.~the electrostatic response of the polarizable continuum.


\emph{This functionality is not yet available for periodic systems.}

%\newpage


\subsection*{Tags for general section of \texttt{control.in}:}

\keydefinition{solvent}{control.in}
{
  \noindent
  Usage: \keyword{solvent} \option{method} \option{cavity} \option{periodicity} \option{d}\\[1.0ex] 
  Purpose: Specifies the desired implicit solvent model including the method which is used to create the cavity. \\[1.0ex]
  \option{method} is a string which specifies the implicit solvent method. Currently, only ``mpe'' is supported. \\[1.0ex]
  \option{cavity} is a string which specifies the method used to create the cavity. Currently, only the isodensity surface cavity (``isodens'') based on the evaluation of the density on a Cartesian grid is supported. \\[1.0ex]
  \option{periodicity} is an integer which specifies the periodicity of the cavity. Currently, only cluster calculations (periodicity 0) are supported. \\[1.0ex]
  \option{d} is a real number (unit unknown). This keyword defines the distance of grid points in the Cartesian grid. Usually, a value of approximately $0.08$ is enough to create a sufficiently large amount of points on the cavity surface (at least 1500). \\
}

\keydefinition{cavity\_surface\_position}{control.in}
{
  \noindent
%solvent                            mpe isodens 0 0.08 
  Usage: \keyword{cavity\_surface\_position} \option{isovalue} \\[1.0ex] 
  Purpose: Specifies the isodensity value for the creation of the isodensity cavity. \\[1.0ex]
  \option{isovalue} is a real number (unit unknown) and specifies the isodensity value. For small organic molecules a value of approximately $0.006$ usually yields reasonable results. Be aware that this number can have a drastic influence on the calculated results. Thus, it does not have a default value but has to be provided by the user when the keyword \keyword{solvent} is present. \\
}

\keydefinition{solvent\_permittivity}{control.in}
{
  \noindent
%solvent                            mpe isodens 0 0.08 
  Usage: \keyword{solvent\_permittivity} \option{epsilon} \\[1.0ex] 
  Purpose: Specifies the macroscopic dielectric constant of the solvent. \\[1.0ex]
  \option{isovalue} is a real number equal to the macroscopic dielectric constant of the solvent. This keyword does not have a default value and has to be provided by the user when the keyword \keyword{solvent} is present. \\
}

\keydefinition{cavity\_rebuild}{control.in}
{
  \noindent
%solvent                            mpe isodens 0 0.08 
  Usage: \keyword{cavity\_rebuild} \option{n} \\[1.0ex] 
  Purpose: Specifies after how many scf steps the cavity is rebuilt. \\[1.0ex]
  \option{n} is an integer. Default: 10000 \\
}
